\documentclass[12pt]{article}
\usepackage[margin=1in]{geometry}
\usepackage{graphicx}
\usepackage{amsmath,amssymb,amsthm}
\usepackage{booktabs}
\usepackage{hyperref}
\usepackage{setspace}
\onehalfspacing

\title{Stablecoins and Parallel Foreign Exchange Markets}
\author{}
\date{}

\begin{document}
\maketitle

\section{Introduction}

Stablecoins have emerged as a new instrument in parallel foreign-exchange markets, particularly in economies with fixed or heavily managed exchange rates. In such settings, discrepancies between official and parallel exchange rates generate incentives for households and firms to seek alternative means of accessing foreign currency. Traditionally, that demand has been met through fragmented cash-based markets. Stablecoins create a new margin of adjustment. They reduce the cost of acquiring foreign currency, provide a publicly visible market price, and thereby reshape both individual hedging behavior and aggregate coordination.

This paper develops a model of parallel foreign-exchange demand in which stablecoins matter through two distinct channels. First, they reduce the effective cost of accessing foreign currency relative to informal cash-based markets. Second, because stablecoin prices are publicly visible and increasingly informative when the stablecoin market is active, they improve common information about the degree of exchange-rate misalignment. These channels jointly determine foreign-currency demand, crisis dynamics, and welfare.

The key conceptual distinction in the model is between a hedge motive and a run motive for foreign-currency demand. The hedge motive reflects households' desire to preserve purchasing power and maintain access to tradable goods when the official exchange rate is overvalued. The run motive reflects strategic demand that arises when households fear a disorderly adjustment and expect others to rush into foreign currency. Stablecoins affect both motives, but they are especially powerful on the second margin because they improve public information and therefore strengthen coordination. This is the paper's central insight: stablecoins are not merely a cheaper vehicle for accessing foreign currency. They are also endogenous public signals that alter the coordination environment in parallel markets.

This perspective generates four main results. First, foreign-currency demand rises monotonically with exchange-rate overvaluation. This occurs because worsening misalignment directly raises the demand for FX as a hedge, even before panic becomes dominant. Second, stablecoin demand also rises monotonically with overvaluation because stablecoins become an increasingly attractive instrument for satisfying that demand. Third, stablecoins increase financial fragility. By lowering FX access costs and improving public information, they raise aggregate FX demand and make it more synchronized across agents, thereby increasing crisis probability. Fourth, the welfare effects of stablecoins are non-monotonic. When exchange-rate misalignment is modest, the access benefit dominates and stablecoins improve welfare. When misalignment is severe, however, the coordination channel dominates and stablecoins reduce welfare by amplifying crisis risk and crisis-state losses.

The paper contributes to three literatures. First, it contributes to the literature on parallel exchange rates and foreign-currency demand by introducing a new digital instrument into the analysis of unofficial FX markets. Existing work typically studies cash-based parallel markets or black-market premia without modeling how digital foreign-currency access changes both the cost and the informational structure of the market. Second, it contributes to the broader literature on public information and coordination by endogenizing the informativeness of the public signal. In standard coordination environments, public information is exogenous. Here, stablecoin participation itself makes the public signal more informative. Third, the paper contributes to the growing literature on stablecoins and digital money by emphasizing a mechanism that is especially relevant in macroeconomically distorted environments: stablecoins can be privately useful while socially destabilizing.

The model is intentionally parsimonious, but it is structured to match the main economic forces at work. Foreign-currency demand is decomposed into hedge demand and run demand. Stablecoins affect both the effective cost of FX access and the precision of public information. Crisis risk depends on whether aggregate FX demand exceeds effective defense capacity. Welfare is defined over both normal and crisis states so that stablecoins can improve access in normal times while worsening disruption in crisis states. Quantitatively, the model can jointly generate monotone FX demand, interior crisis probabilities in both the cash-only and stablecoin economies, higher fragility under stablecoins, and a welfare reversal as overvaluation rises.

These results have a sharp policy implication. Stablecoins should not be evaluated solely as a technological innovation that lowers transaction costs. In economies with overvalued official exchange rates and fragile macro-financial regimes, stablecoins also alter the informational structure of the foreign-exchange market. As a result, they can improve access while at the same time increasing the likelihood of destabilizing runs. The welfare consequences therefore depend on the underlying state of the economy. This state dependence is essential for designing sensible policy toward stablecoins in distorted FX environments.

The remainder of the paper proceeds as follows. Section 2 presents the model. Section 3 develops the theory and formal propositions. Section 4 quantifies the model and shows how the main mechanisms appear in the simulations. Section 5 concludes.

\section{Model}

This section develops a model of parallel foreign-exchange demand in an economy with an overvalued official exchange rate. The model is motivated by environments in which households and firms face incentives to acquire foreign currency outside official channels and may do so using either traditional cash-based markets or stablecoins.

The central mechanism is that stablecoins affect equilibrium through two distinct channels. First, they reduce the effective cost of accessing foreign currency. Second, because stablecoin prices are publicly visible and increasingly informative when usage is high, they improve common information and thereby strengthen coordination among agents. These two effects jointly shape foreign-currency demand, crisis risk, and welfare.

A key feature of the model is that total foreign-currency demand is decomposed into a hedge motive and a run motive. This distinction is central both analytically and quantitatively. The hedge motive captures the desire to preserve purchasing power when the official exchange rate is overvalued. The run motive captures strategic demand arising from concerns about macro-financial instability and the actions of others. Stablecoins matter for both motives, but especially for the second.

\subsection{Environment}

Consider a static economy indexed by a scalar state variable $\theta \ge 0$, which measures the degree of overvaluation of the official exchange rate. Higher values of $\theta$ correspond to a larger wedge between the official and shadow value of foreign currency, greater excess demand for foreign exchange, and weaker ability of the authorities to defend the peg.

The focus on $\theta \ge 0$ reflects the empirically relevant region in which parallel foreign-exchange markets emerge and stablecoins become useful as an alternative vehicle for foreign-currency access.

The economy is populated by a continuum of households of unit mass. Households seek access to foreign currency because foreign-currency holdings provide protection against the erosion of domestic purchasing power and insure against the possibility of disorderly adjustment. They can access foreign currency through two instruments: cash foreign currency and stablecoins.

\subsection{Information}

Households do not directly observe the true degree of overvaluation $\theta$. Each household receives a private signal and, when stablecoins are available, a public signal from the stablecoin market. Public-signal precision is endogenous and improves with stablecoin demand. Let $\omega_s$ denote the posterior weight placed on the public signal. This weight is increasing in stablecoin adoption.

\subsection{Foreign-Currency Demand}

The model's central object is total foreign-currency demand,
\[
Q^{FX} = Q^{hedge} + Q^{run}.
\]

Hedge demand is given by
\[
Q^{hedge} = \max\{0,\chi_h \theta - \eta_h p^{FX}\},
\]
where $p^{FX}$ is the effective cost of obtaining foreign currency.

Run demand is
\[
Q^{run} = \gamma_r (1+\psi \omega_s)\Lambda\left(\frac{\pi-\pi_0}{\tau_r}\right),
\]
where $\pi$ is crisis probability, $\omega_s$ is the public-signal weight, and $\Lambda$ is the logistic function.

This decomposition is central. The hedge motive ensures that FX demand rises smoothly with overvaluation. The run motive captures the coordination channel through which stablecoins amplify fragility.

\subsection{Allocation Across Cash and Stablecoins}

Households allocate total FX demand across cash and stablecoins. Stablecoin usage is increasing in the relative cost advantage of stablecoins and in the quality of public information. As a result, stablecoins account for a larger share of FX demand when overvaluation is severe and the informational role of the stablecoin market becomes more important.

\subsection{Crisis Risk}

Crisis probability is increasing in aggregate FX demand and decreasing in effective defense capacity. Defense capacity declines as overvaluation rises. This creates a natural feedback loop: higher overvaluation raises hedge demand, stronger public information raises run demand, and both margins increase the probability of crisis.

\subsection{Structural Welfare}

Welfare is defined over two states: a normal state and a crisis state. In normal times, stablecoins improve welfare by lowering the effective cost of accessing foreign currency and improving tradable-goods smoothing. In crisis states, however, stablecoins increase disruption losses by strengthening coordination and increasing run intensity. Expected welfare is therefore the weighted average of normal-state and crisis-state utility, with weights determined by crisis probability.

\section{Theory}

This section derives the model's main implications. The key innovation of the framework is the decomposition of foreign-currency demand into a hedge component and a run component. This allows the model to generate monotone foreign-currency demand, endogenous public information, crisis amplification, and welfare reversal within a single tractable structure.

\subsection{Monotone Foreign-Currency Demand}

Total foreign-currency demand is the sum of hedge demand and run demand. Hedge demand rises with overvaluation and falls with FX acquisition cost. Run demand rises with crisis beliefs and with the weight placed on the public signal. Since the hedge motive is fundamentals-driven, the model delivers weakly increasing FX demand as exchange-rate misalignment worsens.

\subsection{Stablecoin Demand}

Stablecoin demand rises with overvaluation because total FX demand rises and because stablecoins have lower access costs than cash. Stablecoins therefore become a larger share of FX demand as the official exchange rate becomes more distorted.

\subsection{Endogenous Information}

Stablecoin usage improves the precision of public information. As the stablecoin market becomes more active, its price becomes more informative, increasing the weight households place on the public signal. This creates a positive feedback between stablecoin adoption and common information.

\subsection{Crisis Amplification}

Stablecoins increase crisis risk through two channels. First, they lower access costs and therefore raise total FX demand. Second, they improve public information and thereby strengthen coordination. Both effects increase the probability that aggregate demand exceeds effective defense capacity.

\subsection{Welfare Reversal}

Stablecoins improve welfare in moderately overvalued states by lowering the cost of FX access and improving tradable-goods smoothing. But they reduce welfare in sufficiently fragile states because they increase coordinated runs and amplify disruption losses in crisis states. The model therefore implies a threshold level of overvaluation below which stablecoins are beneficial and above which they are harmful.

\section{Theory: Formal Results}

This section states the model's main implications in proposition form.

\paragraph{Proposition 1. Monotone foreign-currency demand.}
Suppose the cost schedule is not so steep that the increase in effective FX cost dominates the increase in overvaluation. Then hedge demand is weakly increasing in overvaluation. If crisis probability is also weakly increasing in overvaluation, then total foreign-currency demand is weakly increasing in overvaluation.

\emph{Intuition.} The hedge motive ensures that worsening misalignment directly raises desired FX holdings even absent panic. The run motive then reinforces this increase once crisis beliefs begin to rise.

\paragraph{Proposition 2. Stablecoin demand rises with overvaluation.}
If stablecoins have lower access costs than cash and the stablecoin premium does not rise too quickly with quantity, then stablecoin demand is weakly increasing in overvaluation.

\emph{Intuition.} As total foreign-currency demand rises with overvaluation, the lower-cost instrument captures a growing share of that demand.

\paragraph{Proposition 3. Endogenous public information.}
Public-signal precision is increasing in stablecoin demand. As a result, the weight placed on the public signal is increasing in stablecoin demand.

\emph{Intuition.} Stablecoin adoption does not merely respond to information; it produces information by making the stablecoin price more informative.

\paragraph{Proposition 4. Information-participation feedback.}
There is a positive feedback between stablecoin participation and public information. Greater stablecoin use improves public-signal precision, which increases the weight placed on public information and thereby raises the coordination component of FX demand.

\paragraph{Proposition 5. Stablecoins increase fragility.}
Holding fixed the underlying overvaluation state, the presence of stablecoins weakly increases crisis probability whenever stablecoins reduce FX access costs and improve public information.

\emph{Intuition.} Stablecoins raise total FX demand both by lowering the effective cost of hedging and by increasing the run component through stronger common information.

\paragraph{Proposition 6. Welfare reversal.}
There exists a threshold level of overvaluation such that stablecoins improve welfare below the threshold but reduce welfare above it.

\emph{Intuition.} At low levels of misalignment, the access benefit dominates. At high levels of misalignment, stablecoins amplify coordinated runs and increase expected disruption losses.

\section{Quantitative Results}

This section evaluates the model quantitatively and connects each theoretical mechanism to simulated outcomes.

\subsection{Calibration and Setup}

The model is simulated across a grid of overvaluation levels. Parameters are chosen to ensure three properties: foreign-currency demand increases monotonically in overvaluation, both crisis and non-crisis outcomes occur with positive probability, and stablecoins affect both access costs and information precision.

\subsection{Foreign-Currency Demand}

Foreign-currency demand increases monotonically with overvaluation. The increase is driven initially by hedge demand and later by run demand as crisis beliefs rise. Stablecoin usage rises with overvaluation and accounts for a growing share of total FX demand.

\subsection{Crisis Probability}

Crisis probability increases with overvaluation in both the cash-only and stablecoin economies. However, the presence of stablecoins shifts the crisis probability upward at all levels of overvaluation. This reflects both increased demand and stronger coordination.

\subsection{Welfare}

Welfare exhibits a non-monotonic pattern. In low and moderate overvaluation states, stablecoins improve welfare by lowering access costs and increasing consumption smoothing. In high-overvaluation states, stablecoins reduce welfare by increasing crisis probability and amplifying losses in crisis states.

\subsection{Decomposition}

Counterfactual simulations shutting down individual channels show that both the cost-reduction channel and the information channel are necessary. Removing the cost channel eliminates welfare gains, while removing the information channel eliminates crisis amplification.

\section{Conclusion}

Stablecoins introduce a new channel into parallel foreign-exchange markets by simultaneously lowering FX access costs and improving public information. This dual role makes them privately useful but potentially socially destabilizing. The model shows that stablecoins can improve welfare when exchange-rate misalignment is moderate but reduce welfare when misalignment is severe. The key policy lesson is that evaluating stablecoins in distorted FX environments requires accounting not only for financial access but also for the coordination externalities they generate.

\end{document}
